\documentclass[11pt]{article}

\usepackage[margin=0.9in]{geometry}
\usepackage{amsmath,amssymb,amsthm,mathtools,bm,booktabs,enumitem,microtype}
\usepackage{xcolor}
\usepackage[colorlinks=true,linkcolor=blue!50!black,urlcolor=blue!50!black,citecolor=blue!50!black]{hyperref}

\theoremstyle{plain}
\newtheorem{theorem}{Theorem}
\newtheorem{lemma}[theorem]{Lemma}
\newtheorem{proposition}[theorem]{Proposition}
\newtheorem{corollary}[theorem]{Corollary}

\theoremstyle{definition}
\newtheorem{definition}[theorem]{Definition}

\theoremstyle{remark}
\newtheorem{remark}{Remark}

\DeclareMathOperator{\tr}{tr}
\DeclareMathOperator{\diag}{diag}
\DeclareMathOperator{\rank}{rank}
\DeclareMathOperator{\aut}{aut}
\DeclareMathOperator{\supp}{supp}
\DeclareMathOperator{\He}{He}
\DeclareMathOperator{\RIE}{RIE}
\DeclareMathOperator{\sech}{sech}
\DeclareMathOperator{\pv}{p.v.}
\DeclareMathOperator{\divop}{div}

\newcommand{\E}{\mathbb{E}}
\newcommand{\Prob}{\mathbb{P}}
\newcommand{\R}{\mathbb{R}}
\newcommand{\sL}{\mathsf{L}}
\newcommand{\sS}{\mathsf{S}}
\newcommand{\sR}{\mathsf{R}}
\newcommand{\sZ}{\mathsf{Z}}

\title{\bfseries Resolution of Gaussian Equivalence for Extensive-Rank Matrix Denoising}
\author{}
\date{\vspace{-1.6em}%
\small
\begin{tabular}{@{}c@{}}
Target paper: Z.~He, Q.~Chen, J.~Ma, H.~Zhang, and X.~Yuan,\\
\emph{Extensive-Rank Matrix Denoising Under i.i.d.\ Non-Gaussian Noise},\\
\href{https://arxiv.org/abs/2611.29179}{arXiv:2611.29179}, Section~III-C, Conjecture~1 (\texttt{Conj:Effective})\\[0.35em]
Companion verification script: \texttt{solutions/extensive\_matrix\_denoise/explore.py}\\[0.35em]
November 2026
\end{tabular}}

\begin{document}

\maketitle

\begin{abstract}
We resolve Conjecture~1 (\texttt{Conj:Effective}) of He, Chen, Ma, Zhang, and Yuan (\href{https://arxiv.org/abs/2611.29179}{arXiv:2611.29179}) on orthogonal approximate message passing (OAMP) for extensive-rank matrix denoising under i.i.d.\ non-Gaussian noise.
First, using the bounded continuous extension of the rotationally invariant estimator (RIE) from Appendix~A of the paper together with the \emph{spectral-integral} state evolution (SE) recursion, we prove both parts of Conjecture~1 in probability at every fixed iteration $t \ge 1$ with positive population input variances.
In the independent bi-orthogonally invariant signal model with a compactly supported limiting spectral law $\nu_L$ of positive first moment $p_{\sL} > 0$, the sharp signal hypothesis is convergence of normalized Frobenius power $\|L\|_F^2/(mn) \to p_{\sL}$ in probability (implied by Section~II's almost-sure operator-norm bound, and necessary at $t=1$), while only a finite second moment $0 < p_{\sS} < \infty$ of the entrywise noise prior is required (allowing infinite fourth moment).
The proof develops a finite-graph matrix Hilbert space with simultaneous entrywise Wick--Hermite and spectral noncrossing-tree descriptions, combined with Gaussian integration by parts on smooth odd-polynomial matrix maps $g(X) = X\tilde{g}(X^\top X/n)$.
Second, we disprove unconditional identification of the spectral integral with the limit of expected risks under the paper's almost-sure hypotheses alone: for bounded RIE extensions, uniform integrability of normalized signal power is both necessary and sufficient for their equality, and an explicit full-rank bi-invariant mixture ($\delta = 0.05$) violates part~(ii) at iteration~$1$ (empirical risk $0.449600$ vs.\ expected-risk target $0.499600$) and part~(i) at iteration~$2$ (empirical cross moment $0.128737 \neq 0$).
Finally, for rectangular aspect ratios $\beta < 1$, we establish the exact inverse-square criterion $m^{-1}\sum_{s_i < \eta} s_i^{-2} \to 0$ governing the literal unregularized Hilbert pole below the spectral gap.
\end{abstract}

\section{Model, Calibration, and Summary of Verdicts}\label{sec:setup}

Consider $Y = L + S \in \R^{m \times n}$ with $m \le n$ and $m/n \to \beta \in (0, 1]$ as $m, n \to \infty$.
The signal $L$ and noise $S$ are independent. The law of $L$ is bi-orthogonally invariant (invariant under left and right multiplication by deterministic orthogonal matrices), while the entries of $S$ are i.i.d.\ from a fixed centered prior $P_{\sS}$ with variance $0 < p_{\sS} := \E[S_{ij}^2] < \infty$.
Section~II of \href{https://arxiv.org/abs/2611.29179}{arXiv:2611.29179} assumes that $\nu_{L,n} := m^{-1}\sum_{i=1}^m \delta_{\lambda_i(LL^\top/n)}$ converges weakly almost surely to a deterministic probability measure $\nu_L$ on $\R_+$ and that $\limsup_{n\to\infty}\|L\|_{\mathrm{op}}/\sqrt{n} < \infty$ almost surely, with signal power $p_{\sL} := \int x\,\nu_L(dx) > 0$.
Write $|X|_n := \|X\|_F / \sqrt{mn}$ and $\langle X, Z\rangle_n := \tr(X^\top Z)/(mn)$. Throughout, convergence of random empirical quantities is in probability unless stated otherwise.

\paragraph{Auxiliary Gaussian channels and denoisers.}
For variance $v > 0$:
\begin{itemize}[leftmargin=2em,itemsep=2pt]
\item \emph{Entrywise module:} Let $\sR_v := \sS + \sqrt{v}\,\sZ$ with independent $\sS \sim P_{\sS}$ and $\sZ \sim \mathcal{N}(0,1)$. Take the Gaussian-convolution version of the posterior mean $F_{\sS}(r, v) := \E[\sS \mid \sR_v = r]$ and its risk $m_{\sS}(v) := \E[(\sS - F_{\sS}(\sR_v, v))^2]$.
\item \emph{Spectral module:} Let $\bm{\sR}_v := L + \sqrt{v}\,G$, where $G \in \R^{m\times n}$ has i.i.d.\ $\mathcal{N}(0,1)$ entries independent of $L$. In the paper's $p_{\sL}$-normalized units $s = \sigma/\sqrt{np_{\sL}}$, let $\mu_v$ be the limiting singular-value law of $\bm{\sR}_v/\sqrt{np_{\sL}}$ with density $\rho_v$, symmetrization $\widehat{\mu}_v$, and principal-value Hilbert transform $\mathcal{H}_v(s) := \pv\int_{\R}(s-x)^{-1}\,d\widehat{\mu}_v(x)$. On $\supp(\mu_v) \cap (0,\infty)$, the optimal shrinker of Troiani et~al.\ (2022) is
\begin{equation}\label{eq:shrinker}
f_v(s) = s - \frac{v(1-\beta)}{p_{\sL} s} - \frac{2\beta v}{p_{\sL}}\mathcal{H}_v(s).
\end{equation}
Following Appendix~A of \href{https://arxiv.org/abs/2611.29179}{arXiv:2611.29179}, fix a bounded continuous extension of $f_v$ to $[0,\infty)$ with $f_v(0) = 0$. Given a thin SVD $R = U\diag(\sigma_1,\dots,\sigma_m)V^\top$ with $s_i = \sigma_i/\sqrt{np_{\sL}}$, define $F_{\sL}(R, v) := \sqrt{np_{\sL}}\,U\diag(f_v(s_1),\dots,f_v(s_m))V^\top$.
\end{itemize}
Define the \emph{spectral-integral} calibration and its \emph{expected-risk} counterpart:
\begin{align}
m_{\sL}^{\mathrm{int}}(v)
&:= v - \frac{v^2}{p_{\sL}}\int_{\R_+}\!\left[\frac{(1-\beta)^2}{s^2}
  + \frac{\pi^2\beta^2}{3}\rho_v(s)^2\right]\mu_v(ds),\label{eq:mL-int}\\
m_{\sL}^{\mathrm{exp}}(v)
&:= \lim_{n\to\infty}\frac{1}{mn}\E\bigl\|L - F_{\sL}(\bm{\sR}_v, v)\bigr\|_F^2.\label{eq:mL-exp}
\end{align}
For $D \in \{\sS, \sL\}$, set $a_D(v) := m_D(v)/v$ and $\Phi_D(v) := m_D(v)/(1 - a_D(v))$. Starting from $L^1 := 0$ and $v_1 := p_{\sL}$, the population state evolution (SE) sets $w_t := \Phi_{\sS}(v_t)$ and $v_{t+1} := \Phi_{\sL}(w_t)$, and the OAMP recursion updates
\begin{equation}\label{eq:oamp}
R_{\sS}^t = Y - L^t,\quad
S^t = \frac{F_{\sS}(R_{\sS}^t, v_t) - a_{\sS}(v_t)R_{\sS}^t}{1 - a_{\sS}(v_t)},\qquad
R_{\sL}^t = Y - S^t,\quad
L^{t+1} = \frac{F_{\sL}(R_{\sL}^t, w_t) - a_{\sL}(w_t)R_{\sL}^t}{1 - a_{\sL}(w_t)}.
\end{equation}

\begin{quote}
\textbf{Conjecture~1 (\texttt{Conj:Effective}, \href{https://arxiv.org/abs/2611.29179}{arXiv:2611.29179}).}
\emph{For each fixed $t \ge 1$ with $v_t, w_t > 0$, as $m, n \to \infty$ with $m/n \to \beta$, in probability:}
\begin{enumerate}[label=\textup{(\roman*)}]
\item \emph{With $\sS^t := (F_{\sS}(\sR_{v_t}, v_t) - a_{\sS}(v_t)\sR_{v_t})/(1 - a_{\sS}(v_t))$,}
\[
W_2\!\left(\frac{1}{mn}\sum_{i,j}\delta_{(S_{ij}, R_{\sS,ij}^t, S_{ij}^t)},\;
\mathcal{L}(\sS, \sR_{v_t}, \sS^t)\right) \to 0.
\]
\item \emph{With $\mu_n(M) := m^{-1}\sum_{i=1}^m \delta_{\sigma_i(M)/\sqrt{n}}$,}
\[
W_2\bigl(\mu_n(R_{\sL}^t),\,\mu_n(\bm{\sR}_{w_t})\bigr) \to 0
\qquad\text{and}\qquad
\bigl|L - F_{\sL}(R_{\sL}^t, w_t)\bigr|_n^2 \to m_{\sL}(w_t).
\]
\end{enumerate}
\end{quote}

\paragraph{Summary of verdicts and scope.}
\begin{itemize}[leftmargin=2em,itemsep=2pt]
\item \textbf{Spectral-integral SE + Appendix~A bounded RIE: Proved} (Theorem~\ref{thm:main}) for every \emph{fixed} iteration $t \ge 1$ and \emph{population} SE parameters, assuming only $0 < p_{\sS} < \infty$ and $\mu_n(L) \xrightarrow{W_2} \nu_L^{1/2}$ in probability. Passing to $t \to \infty$ jointly with $n$ or replacing population parameters by data-estimated (e.g., SURE) parameters is not covered.
\item \textbf{Unconditional expected-risk calibration~\eqref{eq:mL-exp}: Disproved} (Theorem~\ref{thm:ui}), with exact uniform-integrability equivalence and an explicit full-rank bi-invariant counterexample failing at iterations~$1$ and~$2$.
\item \textbf{Unregularized Hilbert pole ($\beta < 1$): Exact criterion proved} (Proposition~\ref{prop:pole}).
\end{itemize}

\section{Proof of Gaussian Equivalence Under Spectral-Integral SE}\label{sec:proof}

\begin{theorem}[Fixed-iteration Gaussian equivalence]\label{thm:main}
Assume $L$ and $S$ are independent, $L$ is bi-orthogonally invariant, $S$ has i.i.d.\ centered entries with $0 < p_{\sS} < \infty$, and $\mu_n(L)$ converges in $W_2$ in probability to the square-root pushforward $\nu_L^{1/2}$ of a compactly supported deterministic law $\nu_L$ with $p_{\sL} = \int x\,\nu_L(dx) > 0$ (equivalently, $\nu_{L,n} \Rightarrow \nu_L$ and $|L|_n^2 \to p_{\sL}$ in probability).
With the Appendix~A bounded continuous RIE extension and the population spectral-integral calibration $m_{\sL} = m_{\sL}^{\mathrm{int}}$ in~\eqref{eq:mL-int}, both conclusions~\textup{(i)} and~\textup{(ii)} of Conjecture~1 hold in probability for every fixed $t \ge 1$ with $v_t, w_t > 0$.
Conversely, $|L|_n^2 \to p_{\sL}$ in probability is necessary for~\textup{(i)} at $t = 1$.
\end{theorem}

The necessity of $|L|_n^2 \to p_{\sL}$ at $t=1$ is immediate: applying the continuous quadratic-growth map $(s,r,e) \mapsto (r-s)^2$ to the $W_2$-convergent empirical measure in~(i) gives $|R_{\sS}^1 - S|_n^2 = |L|_n^2 \to \E[(\sR_{v_1} - \sS)^2] = v_1 = p_{\sL}$ in probability. (Likewise $P_{\sS} \in \mathcal{P}_2$ is required for the joint $W_2$ statement.)
We prove sufficiency in five steps.

\subsection{Graph-indexed matrix Hilbert space}

Let $(A_{ij}, S_{ij})_{1\le i\le m,\,1\le j\le n}$ be an i.i.d.\ array where $A_{ij} \sim \mathcal{N}(0,1)$ is independent of $S_{ij} \sim P_{\sS}$.
Choose a countable real orthonormal basis $\mathcal{B}$ of centered functions in $L^2(A, S)$ such that every $\psi \in \mathcal{B}$ has all moments finite and the coordinate function $\psi(A,S)=A$ belongs to $\mathcal{B}$.
For instance, take products of normalized Hermite polynomials in $A$ with a bounded orthonormal basis of $L^2(P_{\sS})$ containing $1$ (constructed by Gram--Schmidt on simple functions from a countable generating algebra), omitting the constant product $1$. Note that $S \in L^2(P_{\sS})$ lies in the closed span of $\mathcal{B}$.

Consider finite simple bipartite graphs $H$ with two distinguished terminals (a left terminal $u_L$ and a right terminal $u_R$) and edges labeled by elements of $\mathcal{B}$, generated recursively from single labeled edges:
\begin{itemize}[leftmargin=2em,itemsep=2pt]
\item A \emph{string} joins an odd number $k \ge 3$ of blocks end-to-end, reversing terminal orientation on every second block, with all internal vertices distinct.
\item A \emph{block} is either a single labeled edge, or at least two parallel branches connecting $u_L$ and $u_R$ with pairwise disjoint interiors, where each branch is a string except that at most one branch may be a single edge.
\end{itemize}
Only finite constructions are used. A string has no direct edge $(u_L, u_R)$, and its interior $H \setminus \{u_L, u_R\}$ is connected. Every multi-branch block is $2$-vertex-connected (every vertex lies on a simple $u_L$--$u_R$ path, and deleting any interior vertex leaves another branch connecting $u_L$ to $u_R$), and deleting either terminal leaves a connected graph. Consequently both decompositions are unique: cutting a string at its articulation vertices recovers its ordered blocks, and deleting $\{u_L, u_R\}$ from a parallel block recovers its string branches together with any direct terminal edge.

Work modulo graph isomorphisms preserving terminals, vertex colors, and edge labels. For such a class $H$ with $l$ internal left vertices and $r$ internal right vertices, define the $m \times n$ matrix
\begin{equation}\label{eq:XH}
X_H(i, j) := \frac{1}{m^{l/2} n^{r/2}\sqrt{\aut(H)}}
\sum_{\phi} \prod_{e \in E(H)} \psi_e\bigl(A_{\phi(e)}, S_{\phi(e)}\bigr),
\end{equation}
where $\phi$ runs over color-preserving maps fixing $(u_L, u_R) \mapsto (i, j)$ that are injective within each vertex color.
At any fixed entry $(i,j)$, distinct classes $H \neq H'$ satisfy $\E[X_H(i,j) X_{H'}(i,j)] = 0$, and $\E[X_H(i,j)^2] \to 1$: for a nonzero second-moment contribution, the two embedded edge sets must coincide with matching edge labels, and there are exactly $\aut(H)$ matching embeddings.

Let $\mathcal{H}_{\mathrm{gr}}$ be the abstract real Hilbert space with orthonormal basis $\{H\}$. We say a sequence of random matrices $M_n \in \R^{m\times n}$ \emph{represents} $h \in \mathcal{H}_{\mathrm{gr}}$ if for finite basis combinations $h_j \to h$ in $\mathcal{H}_{\mathrm{gr}}$,
\[
\lim_{j\to\infty}\limsup_{n\to\infty}\Prob\bigl(|M_n - X_{h_j}|_n > \varepsilon\bigr) = 0
\qquad (\forall \varepsilon > 0).
\]
By Lemma~\ref{lem:entrywise} below, this notion is independent of the approximating sequence $(h_j)$. Because the construction uses successive limits ($n \to \infty$ at fixed finite graph combinations, then $j \to \infty$), neither graph sizes nor polynomial degrees grow with $n$.

\subsection{Entrywise Gaussian and Wick--Hermite description}

\begin{lemma}[Entrywise limit and empirical isometry]\label{lem:entrywise}
At any fixed entry and empirically in probability, the string matrices $(X_H)_{H \text{ string}}$ converge to jointly independent $\mathcal{N}(0,1)$ random variables $(Z_H)_H$, independent of $(A, S)$.
A parallel block with direct terminal edge label $\psi$ (or $1$ if absent) and $k_H$ parallel copies of each string type $H$ represents
\[
\psi(A, S) \prod_H \frac{\He_{k_H}(Z_H)}{\sqrt{k_H!}},
\]
where $\He_k$ are the probabilists' Hermite polynomials.
In particular, $\mathcal{H}_{\mathrm{gr}}$ is empirically isometric to the centered subspace of $L^2(A, S, (Z_H)_H)$, and the closed subspace $\mathcal{G} \subset \mathcal{H}_{\mathrm{gr}}$ spanned by the raw edge $A$ and all strings consists of centered Gaussian coordinates jointly Gaussian with $A$ and independent of $S$.
\end{lemma}

\begin{proof}
Consider a finite product of string entries $X_{H_1}(i,j) \cdots X_{H_p}(i,j)$ whose internal vertex occurrences total $d$; the normalization has order $n^{-d/2}$.
In any term with nonzero expectation, every distinct internal vertex label must appear in at least two factors, because every incident edge must be repeated (basis labels are centered) and each individual embedding is injective.
Terms with strictly fewer than $d/2$ distinct internal vertices vanish as $n \to \infty$.
At equality ($d/2$ distinct internal vertices), each internal vertex has a unique partner in another factor. Across any internal edge of a factor, both endpoints must match into the same partner factor so that the edge is repeated. Connectedness of each string interior propagates this identification across the entire string, including the edges incident to the terminals, yielding a pairing of isomorphic strings with matching edge labels and no further interior identifications. Each such pairing contributes $1$ after normalization (the $\aut(H)$ identifications cancel the normalization), reproducing the Gaussian moment formula for independent standard normals $(Z_H)$. Because strings have no direct edge $(u_L, u_R)$, their limits are independent of $(A_{ij}, S_{ij})$.

For two entries $(i,j)$ and $(i',j')$ with $i \neq i'$ and $j \neq j'$, excluding the other entry's terminals from the injective sums changes either string entry by $O(n^{-1})$ in squared $L^2$ norm (by the identical-embedding second-moment count), after which no string can pair across the two entries; hence the two entries are asymptotically independent. Row- or column-sharing pairs comprise an $O(n^{-1})$ fraction of all pairs, yielding concentration of bounded empirical test functions in probability, while higher even moments provide uniform integrability for polynomial tests.

For a parallel block $B$ containing $k_H$ copies of string type $H$ (and direct edge label $\psi$, or $1$), expand the normalized mean-square error between $X_B$ and $\psi(A,S)\prod_H \He_{k_H}(X_H)/\sqrt{k_H!}$. In mixed moments with $X_B$, injectivity forbids contractions between distinct branches inside the same copy of $B$, leaving precisely the Wick pairings matching the Hermite polynomials; the extra automorphism factor $\prod_H k_H!$ permuting identical branches cancels the $\sqrt{k_H!}$ factors, and the direct edge multiplies as an ordinary function of $(A, S)$. This establishes the mean-square approximation and empirical isometry on all blocks (including single edges), spanning every nonconstant product-basis element of $L^2(A, S, (Z_H)_H)$. Passing from finite combinations to all of $\mathcal{H}_{\mathrm{gr}}$ by $|\cdot|_n$-approximation yields empirical weak and second-moment convergence throughout $\mathcal{H}_{\mathrm{gr}}$.
\end{proof}

\subsection{Spectral trace moments and odd open-word expansion}

\begin{lemma}[Spectral Gaussian equivalence and open words]\label{lem:spectral}
For any finite list of distinct blocks $B_1, \dots, B_p$, the rescaled matrices $X_{B_1}/\sqrt{n}, \dots, X_{B_p}/\sqrt{n}$ have the same limiting normalized mixed trace moments (for dimensionally compatible words with transposes, in expectation and in probability) as independent $m \times n$ i.i.d.\ $\mathcal{N}(0,1)$ matrices divided by $\sqrt{n}$.
Moreover, every odd-length open word $W_n$ of length $k$ in these block matrices (scaled by $n^{-(k-1)/2}$) is represented in $\mathcal{H}_{\mathrm{gr}}$ by a finite linear combination of length-$1$ blocks and strings of these blocks, whose coefficients on length-$1$ blocks equal the corresponding Gaussian mixed inner products.
\end{lemma}

\begin{proof}
Expand a normalized cyclic trace $n^{-1}\tr(M_1 \cdots M_q)$ of $q$ factors $M_r \in \{X_{B_a}/\sqrt{n}, X_{B_a}^\top/\sqrt{n}\}$. For a fixed identification pattern whose union graph $U$ has $V$ distinct vertices, where factor $r$ has $v_r$ vertices (including its two terminals), the normalization and summation yield order
\[
n^{V - 1 - \frac{1}{2}\sum_{r=1}^q (v_r - 1)}
\]
(up to bounded powers of $m/n \to \beta$). For a nonzero expectation, every edge of $U$ must be covered by at least two factors. Let $T$ be any spanning tree of the connected union $U$. Since the restriction of $T$ to factor $r$ is a forest on $v_r$ vertices, factor $r$ covers at most $v_r - 1$ edges of $T$, giving $2(V - 1) \le \sum_{r=1}^q (v_r - 1)$. Equality requires that for \emph{every} spanning tree $T$, every edge of $T$ is covered exactly twice and every factor restricts to a spanning tree of its own vertices. Therefore:
\begin{enumerate}[label=\textup{(\arabic*)},leftmargin=2em,itemsep=1pt]
\item Every edge of $U$ is covered exactly twice.
\item No path in $U$ using only edges outside factor $r$ can connect two distinct vertices of factor $r$ with interior disjoint from factor $r$ (otherwise extending that path to a spanning tree $T$ would disconnect $T$ restricted to factor $r$).
\end{enumerate}
Hence each factor is a complete biconnected component of $U$ (or a bridge when the block is a single edge), traversed by exactly two factors sharing the same edge set. Those two factors traverse the component between the same pair of terminals in opposite directions, and all other components attach only at those two terminals. Collapsing each paired block to one edge produces a double-walked tree—identical to the leading diagrams of independent rectangular Gaussian matrices. Conversely, decorating each edge of such a tree with its paired block and fresh internal vertices cancels $\aut(B)$ via root-preserving isomorphisms and cancels all internal vertex dimensions against the block normalizations, preserving exact powers of $\beta$. For the trace variance, diagrams connecting the two trace cycles lose at least one extra power of $n$ by the same spanning-tree bound, while disconnected diagrams factor asymptotically; this yields convergence in probability.

For an odd open word of length $k$ scaled by $n^{-(k-1)/2}$, define a \emph{restricted word} by restricting the sum to embeddings injective on the entire string of blocks (except at the consecutive terminal joins). For $k = 1$ this is the block itself; for $k \ge 3$ it equals $(m/n)^{(k-1)/4}$ times the corresponding string matrix (half the $k-1$ connecting vertices are left vertices, and $\aut$ multiplies across ordered blocks). Form the finite linear combination of restricted words obtained by taking all noncrossing partial pairings of positions in the word such that paired blocks have matching labels, opposite traversal directions, and no unpaired position inside a pair, weighting each deleted pair by $d/n$ where $d \in \{m, n\}$ is the dimension at the \emph{end} of its first traversal. To prove that the difference from the unrestricted word has vanishing mean squared $|\cdot|_n$-norm, close one open word against the reverse of another and apply the cyclic tree count (with any required injectivity constraints). Every leading closed tree splits at the two join vertices into two half-walks sharing the unique simple path between those joins (traversed in reverse in the second half), together with noncrossing closed excursions within each half; each excursion introduces one fresh vertex of dimension $d$ contributing weight $d/n$. An already restricted half has no excursions. Thus the three inner-product terms in the expanded squared difference have identical limits with coefficients $1, -2, 1$. The coefficients on length-$1$ blocks therefore match the Gaussian mixed inner products.
\end{proof}

\subsection{Signal coupling, spectral transfer, and entrywise continuity}

\paragraph{Signal coupling.}
Couple $L$ to the thin SVD of $A$: give $L$ the same left and right singular vectors as $A$ (using independent Haar factors to handle redundant signs) and insert an independent copy of the ordered singular values of $L$, keeping $S$ independent. By bi-orthogonality, $(L, S)$ has the prescribed joint law. Let $h$ be the quantile map from the continuous limiting singular law of $A/\sqrt{n}$ to $\nu_L^{1/2}$ (with bounded endpoint values). Because the singular values are ordered with matching weak limit and convergent second moments, $|L - h(A)|_n^2 = W_2^2(\mu_n(L), \mu_n(h(A))) \to 0$ in probability. Approximating $h$ by odd polynomials in $L^2$ of the Gaussian singular law (which has no atom at zero and compact support) and applying Lemma~\ref{lem:spectral} places $L \in \mathcal{G}$ with $|L|_n^2 \to p_{\sL}$. In particular, $L$ is entrywise Gaussian even when $\nu_L$ has atoms.

\begin{lemma}[Spectral transfer]\label{lem:spectral-transfer}
Suppose $N \in \mathcal{H}_{\mathrm{gr}}$ lies in the closed subspace spanned by blocks orthogonal to $A$, with $|N|_n^2 \to w > 0$. In a comparison model replacing each block by an independent $m \times n$ i.i.d.\ $\mathcal{N}(0,1)$ matrix (designating the $A$-matrix as signal seed $A_{\mathrm{c}}$), the comparison noise $N_{\mathrm{c}}$ is i.i.d.\ $\mathcal{N}(0, w)$ independent of $L_{\mathrm{c}} = h(A_{\mathrm{c}})$. Then:
\begin{enumerate}[label=\textup{(\arabic*)},leftmargin=2em,itemsep=1pt]
\item $W_2(\mu_n(L + N),\, \mu_n(L + \sqrt{w}\,G)) \to 0$ in probability, with the deterministic Gaussian-channel limit $\mu$.
\item $F_{\sL}(L + N, w)$ represents a vector in $\mathcal{H}_{\mathrm{gr}}$, and its limiting $|\cdot|_n$-norm as well as its limiting inner products with $L$ and with every block agree with those of $F_{\sL}(L_{\mathrm{c}} + N_{\mathrm{c}}, w)$ in the Gaussian comparison.
\end{enumerate}
\end{lemma}

\begin{proof}
Approximate $L$ by a fixed odd polynomial in $A$ and $N$ by a fixed finite sum of blocks orthogonal to $A$. By Lemma~\ref{lem:spectral}, all mixed trace moments of the approximating input agree with its Gaussian comparison, whose limiting singular law is compactly supported (finitely many normalized Gaussian matrices have almost surely eventually bounded operator norms). Higher-moment control and compact support give $W_2$ convergence to the same deterministic limit; returning to $(L, N)$ and $(L_{\mathrm{c}}, N_{\mathrm{c}})$ costs at most the $|\cdot|_n$ approximation error by the singular-value $W_2$ perturbation inequality. Since $L_{\mathrm{c}}$ and $L$ differ by $o_P(1)$ in $|\cdot|_n$ under coupling, the triangle inequality through the deterministic limit yields~(1).

For~(2), extend the bounded continuous singular profile of Appendix~A oddly to $\R$ and approximate it uniformly by $K$-Lipschitz odd profiles. An odd $K$-Lipschitz singular map is $K$-Lipschitz in Frobenius norm (apply the scalar function to the symmetric dilation and expand in the two eigenbases, bounding each divided difference by $K$). Thus input approximations pass through $F_{\sL}$ in $|\cdot|_n$. For each fixed approximating input and Lipschitz profile, approximate the profile on the compact limiting support by odd polynomials (with vanishing tail error by higher-moment convergence). Expanding these odd polynomials into odd open words and applying Lemma~\ref{lem:spectral} places them in $\mathcal{H}_{\mathrm{gr}}$ and transfers all required inner products. By empirical isometry (Lemma~\ref{lem:entrywise}), the approximations are Cauchy in $\mathcal{H}_{\mathrm{gr}}$, proving~(2).
\end{proof}

\begin{lemma}[Entrywise continuity under finite variance]\label{lem:entry-cont}
For any prior $P_{\sS}$ with $0 < p_{\sS} < \infty$ and any $v > 0$, the posterior mean $r \mapsto F_{\sS}(r, v)$ is continuous with at most linear growth. Consequently, if $R_{\sS}$ represents a vector in $\mathcal{H}_{\mathrm{gr}}$, then $F_{\sS}(R_{\sS}, v)$ preserves joint empirical $W_2$ convergence and its centered part represents a vector in $\mathcal{H}_{\mathrm{gr}}$.
\end{lemma}

\begin{proof}
Continuity follows from Gaussian convolution. Pick $K > 0$ such that $c := P_{\sS}([-K, K]) > 0$. For any $r \in \R$ and $u \ge 0$, every $s \in [-K, K]$ satisfies $|r - s| \le |r| + K$, whereas every $|s| > 2|r| + K + u$ satisfies $|r - s| > |r| + K + u$; hence
\[
\frac{\exp(-(r-s)^2/(2v))}{\int_{-K}^{K} \exp(-(r-t)^2/(2v))\,P_{\sS}(dt)}
\le \frac{1}{c}\exp\!\left(-\frac{u^2}{2v}\right),
\]
giving $\Prob(|\sS| > 2|r| + K + u \mid \sR_v = r) \le c^{-1} e^{-u^2/(2v)}$. Integrating over $u \ge 0$ yields $|F_{\sS}(r, v)| \le 2|r| + K + c^{-1}\sqrt{\pi v / 2}$. Continuous maps of at most linear growth preserve joint empirical $W_2$ convergence (by truncation and uniform integrability of squared coordinates). Approximating the desired centered output in $L^2$ of the limiting coordinate law by finite product-basis sums (Lemma~\ref{lem:entrywise}) transfers the $L^2$ error to $|\cdot|_n$ via joint $W_2$ convergence, placing the centered output in $\mathcal{H}_{\mathrm{gr}}$.
\end{proof}

\subsection{Induction and smooth odd-polynomial Stein/IBP calibration}

\begin{proof}[Proof of Theorem~\ref{thm:main}]
We prove by induction on $t \ge 1$ that $E^t := L - L^t \in \mathcal{G}$ with $|E^t|_n^2 \to v_t$, together with parts~(i) and~(ii) at iteration~$t$. For $t = 1$, $E^1 = L \in \mathcal{G}$ and $|E^1|_n^2 \to p_{\sL} = v_1$.

\emph{Entrywise step.} Suppose $E^t \in \mathcal{G}$ with $|E^t|_n^2 \to v_t > 0$. By Lemma~\ref{lem:entrywise}, $(S, E^t)$ has empirical $W_2$ limit $(\sS, \sqrt{v_t}\sZ)$ with $\sZ \perp \sS$. Lemma~\ref{lem:entry-cont} applied to $R_{\sS}^t = S + E^t$ and its extrinsic update yields part~(i) of Conjecture~1 and places $N^t := S - S^t$ in $\mathcal{H}_{\mathrm{gr}}$. Set $M := m_{\sS}(v_t)$ and $a := M/v_t$. Posterior orthogonality gives $\E[F_{\sS}(\sR_{v_t}, v_t)\,\sR_{v_t}] = p_{\sS}$ and $\E[F_{\sS}(\sR_{v_t}, v_t)\,\sS] = p_{\sS} - M$, hence $\E[F_{\sS}(\sR_{v_t}, v_t)\,\sqrt{v_t}\sZ] = M$. Therefore $N^t$ is centered, orthogonal to $E^t$, and satisfies $|N^t|_n^2 \to M/(1 - a) = w_t$. Furthermore, $N^t$ is orthogonal to $A$ and to every string: regressing any $Z \in \mathcal{G}$ on $E^t$ leaves an independent Gaussian remainder jointly independent of $(\sS, \sR_{v_t})$. By Lemma~\ref{lem:entrywise}, the remaining components of $N^t$ lie in the closed subspace of blocks orthogonal to $A$.

\emph{Gaussian-reference bulk Gram identities via smooth odd-polynomial IBP.}
Consider the canonical bounded comparison $X = L_{\mathrm{c}} + N_{\mathrm{c}}$ of Lemma~\ref{lem:spectral-transfer} with $N_{\mathrm{c}} = \sqrt{w}\,G$, $w = w_t > 0$. Work in native $\sqrt{n}$ units, denoting the limiting singular law of $X/\sqrt{n}$ by $\mu$ and its optimal on-support shrinker by $f(s) = s - w(1-\beta)/s - 2\beta w\,\mathcal{H}(s)$.
For any odd polynomial $g(s) = s\,\tilde{g}(s^2)$, the rescaled singular-value map
\[
g(X) := X\,\tilde{g}\!\left(\frac{X^\top X}{n}\right)
\]
is a \emph{polynomial—and therefore $C^\infty$ smooth—function of the matrix entries of $X$}, completely avoiding SVD differentiability issues at repeated or zero singular values.
By Gaussian integration by parts (Stein's lemma) on $N_{\mathrm{c}} = \sqrt{w}\,G$, $\E\langle g(X), N_{\mathrm{c}}\rangle_n = \frac{w}{mn}\E[\divop_X g(X)]$.
On the dense open set of matrices with simple positive singular values $\sigma_i = \sqrt{n}s_i$, the SVD divergence identity gives
\[
\frac{1}{mn}\divop_X g(X)
= \frac{n-m}{n}\cdot\frac{1}{m}\sum_{i=1}^m \frac{g(s_i)}{s_i}
  + \frac{m}{n}\cdot\frac{1}{m^2}\sum_{i \neq j}\frac{s_i g(s_i) - s_j g(s_j)}{s_i^2 - s_j^2}
  + \frac{1}{n}\cdot\frac{1}{m}\sum_{i=1}^m g'(s_i),
\]
which extends continuously to all $X \in \R^{m\times n}$ because $g(X)$ is smooth and the diagonal divided difference $(sg(s) - ug(u))/(s^2 - u^2)$ is removable for odd polynomials. As $m/n \to \beta$, the derivative term is $O(1/n) \to 0$. Because $L_{\mathrm{c}} = h(A_{\mathrm{c}})$ has bounded operator norm $\|L_{\mathrm{c}}\|_{\mathrm{op}}/\sqrt{n}$ and $G/\sqrt{n}$ has sub-Gaussian operator-norm tails, polynomial approximation in $A_{\mathrm{c}}$ and trace concentration (Lemma~\ref{lem:spectral}) upgrade the expectation limit to convergence in probability:
\begin{align*}
\langle g(X), N_{\mathrm{c}}\rangle_n
&\longrightarrow w\left[(1-\beta)\int_{\R_+} \frac{g(s)}{s}\,d\mu(s)
  + \beta\iint_{\R_+^2} \frac{s g(s) - u g(u)}{s^2 - u^2}\,d\mu(s)\,d\mu(u)\right] \\
&= \int_{\R_+} g(s)\bigl(s - f(s)\bigr)\,d\mu(s).
\end{align*}
Here the last equality symmetrizes the double integral after truncating to $|s - u| > \varepsilon$: since $\frac{sg(s)-ug(u)}{s^2-u^2}$ is symmetric in $(s,u)$ and $\frac{s}{s^2-u^2} = \frac{1}{2}(\frac{1}{s-u} + \frac{1}{s+u})$, the truncated double integral equals $2\int_{\R_+} g(s)\bigl(\int_{|s-u|>\varepsilon}\frac{s}{s^2-u^2}\,d\mu(u)\bigr)\,d\mu(s)$, which converges to $2\int_{\R_+} g(s)\mathcal{H}(s)\,d\mu(s)$ as $\varepsilon \downarrow 0$ by $L^2$ boundedness of the Hilbert transform on the compactly supported bounded density of $\mu$.
Subtracting from $\langle g(X), X\rangle_n \to \int s g(s)\,d\mu(s)$ gives $\langle g(X), L_{\mathrm{c}}\rangle_n \to \int g(s)f(s)\,d\mu(s)$.
Taking $g(s) = s$ and approximating $g \to f$ yields (with $F_{\mathrm{c}} := F_{\sL}(X, w)$)
\begin{equation}\label{eq:gram}
\langle F_{\mathrm{c}}, X\rangle_n \to p_{\sL},
\qquad
\langle F_{\mathrm{c}}, L_{\mathrm{c}}\rangle_n \to p_{\sL} - m_{\sL}^{\mathrm{int}}(w),
\qquad
|F_{\mathrm{c}}|_n^2 \to p_{\sL} - m_{\sL}^{\mathrm{int}}(w),
\end{equation}
using $m_{\sL}^{\mathrm{int}}(w) = p_{\sL} - \int f(s)^2\,d\mu(s)$ from Appendix~A of \href{https://arxiv.org/abs/2611.29179}{arXiv:2611.29179}.

\emph{Spectral transfer and closing the induction.}
Applying Lemma~\ref{lem:spectral-transfer} to $R_{\sL}^t = L + N^t$ transfers~\eqref{eq:gram} to $F_{\sL}(R_{\sL}^t, w_t)$ and proves both assertions of part~(ii).
It follows directly from~\eqref{eq:gram} and $a_{\sL}(w_t) = m_{\sL}^{\mathrm{int}}(w_t)/w_t$ that $L^{t+1}$ is asymptotically orthogonal to $N^t$ and satisfies $|L - L^{t+1}|_n^2 \to \Phi_{\sL}(w_t) = v_{t+1}$.
Moreover, in the Gaussian comparison, any finite linear combination of blocks orthogonal to both $A$ and $N^t$ corresponds to an independent Gaussian matrix $W$ relative to $(A_{\mathrm{c}}, N_{\mathrm{c}})$, so conditioning on $(A_{\mathrm{c}}, N_{\mathrm{c}})$ gives $\E[\langle L_{\mathrm{c}}^{t+1}, W\rangle_n^2 \mid A_{\mathrm{c}}, N_{\mathrm{c}}] = O(|L_{\mathrm{c}}^{t+1}|_n^2/(mn)) \to 0$.
By Lemma~\ref{lem:spectral-transfer}(2), $L^{t+1}$ is orthogonal in $\mathcal{H}_{\mathrm{gr}}$ to every block perpendicular to $A$. Since $L \in \mathcal{G}$ already has this orthogonality, $E^{t+1} = L - L^{t+1}$ lies in $\mathcal{G}$ with $|E^{t+1}|_n^2 \to v_{t+1}$. This completes the induction.
\end{proof}

\section{Expected-Risk Calibration: Sharp Criterion and Counterexample}\label{sec:ui}

Equation~\eqref{eq:mL-exp} in \href{https://arxiv.org/abs/2611.29179}{arXiv:2611.29179} defines $m_{\sL}(w)$ as the limit of expected risks, equating it to the spectral integral~\eqref{eq:mL-int}. The next result characterizes exactly when the two coincide and disproves their unconditional equality under Section~II's almost-sure hypotheses.

\begin{theorem}[Uniform-integrability equivalence and iteration-$1/2$ counterexample]\label{thm:ui}
Fix $w > 0$, let $Q_n := |L|_n^2$, and let $D_n := |L - F_{\sL}(L + \sqrt{w}\,G, w)|_n^2$ for a bounded normalized singular profile ($|F_{\sL}|_n \le C_w$ deterministically).
\begin{enumerate}[label=\textup{(\arabic*)},leftmargin=2em,itemsep=2pt]
\item Assuming $D_n \to m_{\sL}^{\mathrm{int}}(w)$ in probability,
\[
\E[D_n] \to m_{\sL}^{\mathrm{int}}(w)
\iff (D_n)_{n\ge 1} \text{ is uniformly integrable}
\iff (Q_n)_{n\ge 1} \text{ is uniformly integrable}.
\]
In particular, given $Q_n \to p_{\sL}$ in probability, $\E[D_n] \to m_{\sL}^{\mathrm{int}}(w)$ holds if and only if $\E[Q_n] \to p_{\sL}$.
\item There exists a sequence of full-rank bi-orthogonally invariant signals $L$ with uniformly bounded entrywise second moments satisfying all almost-sure hypotheses of Section~II of \href{https://arxiv.org/abs/2611.29179}{arXiv:2611.29179} for which $m_{\sL}^{\mathrm{exp}}(w_1) \neq m_{\sL}^{\mathrm{int}}(w_1)$, causing the expected-risk calibration to fail part~\textup{(ii)} at iteration~$1$ and part~\textup{(i)} at iteration~$2$.
\end{enumerate}
\end{theorem}

\begin{proof}
\emph{Part~(1).} Since $|F_{\sL}|_n \le C_w$ deterministically, the triangle inequality gives
\[
D_n \le 2Q_n + 2C_w^2, \qquad Q_n \le 2D_n + 2C_w^2.
\]
Hence uniform integrability of $(Q_n)$ is equivalent to that of $(D_n)$. Because $D_n \ge 0$ converges in probability to the finite constant $m_{\sL}^{\mathrm{int}}(w)$, Vitali's convergence theorem gives $\E[D_n] \to m_{\sL}^{\mathrm{int}}(w) \iff (D_n)$ is uniformly integrable.

\emph{Part~(2).} Fix $m = n$, $p_{\sL} = p_{\sS} = 1$, and take i.i.d.\ Rademacher entries $S_{ij} \in \{\pm 1\}$. Independently of $S$, draw $L$ as follows: with probability $1 - n^{-2}$ set $L = A$ (i.i.d.\ $\mathcal{N}(0, 1)$ entries), and with probability $n^{-2}$ set
\[
L = n\sqrt{\delta}\,\sqrt{n}\,O_n, \qquad \delta := 0.05,
\]
where $O_n \sim \operatorname{Haar}(O_n)$. Both components are full rank and bi-orthogonally invariant. Since $\sum_{n\ge 1} n^{-2} < \infty$, the Borel--Cantelli lemma implies that almost surely $L = A$ for all sufficiently large $n$, so Section~II's almost-sure weak convergence and operator-norm bound hold with the standard square Gaussian signal law ($p_{\sL} = 1$). Every entry has uniformly bounded second moment $\E[L_{ij}^2] = (1 - n^{-2}) + \delta \to 1.05$, yet $\E[Q_n] \to 1 + \delta = 1.05 \neq 1$.

Write $M := m_{\sS}(1) = \E[\sech^2(1 + \sZ)]$, $w := M/(1 - M)$, $b := w/(1 + w) \approx 0.449600$, and $c := 1/(1 + w)$. On the Gaussian spectral reference, the on-support shrinker is $cs$; take the Appendix~A bounded continuous extension $f(s) = c\min(s, 2\sqrt{1+w})$, which satisfies $|F_{\sL}|_n \le 2/\sqrt{1+w}$. On the ordinary Gaussian branch, the reference risk converges to $b$ in probability and in expectation (by uniform integrability of $|A|_n^2$). On the exceptional branch, $|L|_n = n\sqrt{\delta}$ while $|F_{\sL}|_n = O(1)$, so its contribution to $\E[D_n]$ is $n^{-2}(n\sqrt{\delta} + O(1))^2 \to \delta$. Thus
\[
m_{\sL}^{\mathrm{exp}}(w) = b + \delta \approx 0.499600 \neq 0.449600 \approx b = m_{\sL}^{\mathrm{int}}(w).
\]
At iteration~$1$ of OAMP on the ordinary branch (which has probability $1 - n^{-2} \to 1$), the entries of $R_{\sL}^1$ are i.i.d.\ copies of $H := (S + A - F_{\sS}(S+A))/(1 - M)$, which are centered with $\E[AH] = 1$, $\E[SH] = w$, and $\E[H^2] = 1 + w$ by scalar posterior orthogonality. By the finite-variance Marchenko--Pastur theorem and the law of large numbers for second moments, $\mu_n(R_{\sL}^1)$ converges in $W_2$ to the scaled quarter-circle law supported on $[0, 2\sqrt{1+w}]$. Hence $|F_{\sL}(R_{\sL}^1) - cR_{\sL}^1|_n \to 0$, and the law of large numbers yields unconditionally in probability
\[
\bigl|L - F_{\sL}(R_{\sL}^1)\bigr|_n^2 \longrightarrow 1 - 2c + c^2(1 + w) = b \approx 0.449600 \neq b + \delta = m_{\sL}^{\mathrm{exp}}(w),
\]
violating part~(ii) of Conjecture~1 at $t = 1$.

If OAMP then applies the expected-risk correction $a' := (b + \delta)/w$, observe that $a' < 1 \iff \delta < M^2/(1 - M)$. Restricting the Gaussian expectation $M = \E[\sech^2(1+\sZ)]$ to $|1+\sZ| \le 1/2$ and $1/2 < |1+\sZ| \le 1$ gives $M > 0.21$, whence $M^2/(1 - M) > 0.0558 > 0.05 = \delta$, so $a' < 1$ and the next SE variance is strictly positive. On the ordinary branch, $L^2 = k R_{\sL}^1 + o_P(1)$ in $|\cdot|_n$ where $k := (c - a')/(1 - a')$, so the normalized empirical cross moment satisfies
\[
\langle S,\, R_{\sS}^2 - S\rangle_n \longrightarrow -kw \approx 0.128737 \neq 0.
\]
Since $(s, r, e) \mapsto s(r - s)$ is a continuous quadratic-growth test under $W_2$ whose expectation under every independent Gaussian channel $\sR_v = \sS + \sqrt{v}\,\sZ$ vanishes, part~(i) of Conjecture~1 fails at $t = 2$. Under the spectral-integral calibration $a_{\sL} = c$, one has $k = 0$ and Theorem~\ref{thm:main} holds.
\end{proof}

\section{Off-Support Extensions and the \texorpdfstring{$\beta < 1$}{beta < 1} Hilbert Pole}\label{sec:pole}

Any bounded continuous extension of~\eqref{eq:shrinker} agreeing on the positive support gives the same in-probability Theorem~\ref{thm:main}. For $\beta = 1$, any continuous extension of at most linear growth also yields the same in-probability limits under spectral-integral calibration by the $W_2$ tail bound (though Theorem~\ref{thm:ui}(1) uses boundedness of the extension).

When $\beta < 1$, continuing the unregularized main-text formula~\eqref{eq:shrinker} literally to all $s > 0$ introduces a pole $-w(1-\beta)/s$ (in native $\sqrt{n}$ units) below the reference spectral gap. The next proposition isolates the exact extra condition required to transfer Theorem~\ref{thm:main} to that unregularized rule.

\begin{proposition}[Literal unregularized Hilbert-pole criterion for $\beta < 1$]\label{prop:pole}
Suppose $\beta < 1$ and an input $R \in \R^{m\times n}$ satisfies $\mu_n(R) \xrightarrow{W_2} \mu$ in probability, where $\mu$ is the Gaussian-reference singular law (in native $\sqrt{n}$ units) supported strictly away from $0$. Fix $\eta > 0$ below the lower edge of $\supp(\mu)$. Let $F_{\sL}^{\mathrm{lit}}(R)$ apply the literal unregularized formula~\eqref{eq:shrinker} to every singular value $s_i = \sigma_i(R)/\sqrt{n}$ (counting $s_i = 0$ as yielding $+\infty$ in the pole), and let $F_{\sL}^{\mathrm{bdd}}(R)$ be the Appendix~A bounded extension. Then:
\begin{enumerate}[label=\textup{(\arabic*)},leftmargin=2em,itemsep=1pt]
\item $|F_{\sL}^{\mathrm{lit}}(R) - F_{\sL}^{\mathrm{bdd}}(R)|_n \to 0$ in probability if and only if
\begin{equation}\label{eq:pole-crit}
\frac{1}{m}\sum_{s_i < \eta} s_i^{-2} \xrightarrow{\Prob} 0.
\end{equation}
\item Under the hypotheses of Theorem~\ref{thm:main}, condition~\eqref{eq:pole-crit} is also necessary for $|L - F_{\sL}^{\mathrm{lit}}(R)|_n^2 \to m_{\sL}^{\mathrm{int}}(w)$ in probability.
\end{enumerate}
\end{proposition}

\begin{proof}
On $[\eta, \infty)$, the literal formula is continuous where evaluated, agrees with $F_{\sL}^{\mathrm{bdd}}$ on $\supp(\mu)$, and grows at most linearly as $s \to \infty$; hence $W_2$ convergence of $\mu_n(R)$ forces the squared per-entry discrepancy from singular values $s_i \ge \eta$ to vanish in probability. On $(0, \eta)$, the discrepancy is asymptotic to $w(1-\beta)/s$, so its squared per-entry norm is $\Theta(m^{-1}\sum_{s_i < \eta} s_i^{-2})$, proving~(1).

For~(2), let $P_\eta$ be the orthogonal projection onto the subspace spanned by left singular vectors of $R$ with $s_i < \eta$. Since $\mu([0, \eta]) = 0$, $\rank(P_\eta) = o_P(m)$. By $W_2$ convergence of $\mu_n(L)$ to the compactly supported law $\nu_L^{1/2}$ and the singular-value variational bound (Ky Fan inequality), $|P_\eta L|_n^2 \to 0$ in probability. Likewise $|P_\eta F_{\sL}^{\mathrm{bdd}}(R)|_n^2 \to 0$ by boundedness of the profile and $\rank(P_\eta) = o_P(m)$. Hence
\[
\bigl|L - F_{\sL}^{\mathrm{lit}}(R)\bigr|_n^2
= \bigl|(I - P_\eta)(L - F_{\sL}^{\mathrm{lit}}(R))\bigr|_n^2
  + \bigl|P_\eta F_{\sL}^{\mathrm{lit}}(R)\bigr|_n^2 + o_P(1).
\]
The first term converges to $m_{\sL}^{\mathrm{int}}(w)$ by the $[\eta,\infty)$ control above, so matching $m_{\sL}^{\mathrm{int}}(w)$ requires $|P_\eta F_{\sL}^{\mathrm{lit}}(R)|_n^2 \to 0$, which is equivalent to~\eqref{eq:pole-crit}.
\end{proof}

Theorem~\ref{thm:main} does not establish automatic control of extreme small singular values for the unregularized pole rule; the Appendix~A bounded extension requires no inverse moments.

\section{Reproducible Verification and Literature Context}\label{sec:verify}

The companion script \texttt{solutions/extensive\_matrix\_denoise/explore.py} (default seed \texttt{261129179}) runs population-calibrated OAMP over three independent replicates at $n \in \{64, 128, 256\}$ for three iterations, testing both the nonlinear Haar orthogonal signal $L = \sqrt{n}\,O_n$ and an i.i.d.\ Gaussian signal $L$, each against four unit-variance entrywise priors (Rademacher, Laplace, Student $t_3/\sqrt{3}$ with infinite fourth moment, and Gaussian).
Table~\ref{tab:haar} reports replicate means at iteration~$3$ for the nonlinear Haar orthogonal signal; the triple coupling column is an explicit upper bound on the empirical three-dimensional $W_2$ distance to an independent finite reference sample. Across all $24$ ordinary model/prior/iteration combinations at $n = 256$, the largest absolute replicate-mean spectral risk gap is $0.002152$.

\begin{table}[ht]
\centering
\caption{Nonlinear Haar orthogonal signal ($L = \sqrt{n}\,O_n$), iteration $t = 3$, replicate means (seed \texttt{261129179}).}
\label{tab:haar}
\smallskip
\begin{tabular}{@{}lrrrr@{}}
\toprule
& \multicolumn{2}{c}{Spectral $W_2$} & Triple coupling & Spectral risk gap \\
\cmidrule(lr){2-3}
Entry prior & $n = 64$ & $n = 256$ & ($n = 256$) & ($n = 256$) \\
\midrule
Rademacher    & 0.04297 & 0.01302 & 0.21935 & $+0.00122$ \\
Laplace       & 0.04525 & 0.00995 & 0.21066 & $+0.00025$ \\
Student $t_3$ & 0.03465 & 0.00895 & 0.66826 & $+0.00065$ \\
Gaussian      & 0.03515 & 0.01094 & 0.14203 & $-0.00082$ \\
\bottomrule
\end{tabular}
\end{table}

Table~\ref{tab:rare} verifies the rare-event mixture counterexample of Theorem~\ref{thm:ui}(2) ($\delta = 0.05$) by stratified Monte Carlo across the two branches with exact mixture weights $1 - n^{-2}$ and $n^{-2}$, together with the iteration-$2$ cross moment under expected-risk calibration.

\begin{table}[ht]
\centering
\caption{Rare-event bi-invariant counterexample ($\delta = 0.05$, Rademacher $S$): empirical vs.\ analytic limits.}
\label{tab:rare}
\smallskip
\begin{tabular}{@{}rrrr@{}}
\toprule
$n$ & Ordinary reference risk & Weighted reference risk & Iteration-$2$ cross moment \\
\midrule
64  & 0.444842 & 0.484913 & 0.137741 \\
128 & 0.448012 & 0.492936 & 0.123486 \\
256 & 0.450177 & 0.497612 & 0.132802 \\
\midrule
Analytic limit & 0.449600 & 0.499600 & 0.128737 \\
\bottomrule
\end{tabular}
\end{table}

\paragraph{Literature context.}
The graph-indexed matrix calculus and spanning-tree / biconnected-component counting in Section~\ref{sec:proof} connect naturally to Male's traffic-freeness framework for permutation-invariant random matrices (\emph{Traffic Distributions and Independence}, \emph{Mem.\ Amer.\ Math.\ Soc.}\ 267, no.~1300, 2020; \href{https://arxiv.org/abs/1111.4662}{arXiv:1111.4662}) and to the tensor-network moment methods for AMP universality and orthogonally invariant ensembles of Wang, Zhong, and Fan (\emph{Ann.\ Appl.\ Probab.}\ 34(4), 2024, \href{https://arxiv.org/abs/2206.13037}{arXiv:2206.13037}; \emph{Inf.\ Inference} 13(3), 2024, \href{https://arxiv.org/abs/2110.02318}{arXiv:2110.02318}).
The Gaussian-channel optimal rectangular RIE and its risk integral are due to Troiani, Erba, Krzakala, Maillard, and Zdeborov\'a (PMLR~190, 2022; \href{https://arxiv.org/abs/2203.07752}{arXiv:2203.07752}), as cited in \href{https://arxiv.org/abs/2611.29179}{arXiv:2611.29179}.
Our proof applies to fixed iterations $t$, known limiting priors, and population SE coefficients; extending Gaussian equivalence to $t \to \infty$ jointly with $n$ or to data-estimated (e.g., SURE) calibration remains open.

\end{document}
